\subsection{Discussion of the research questions}
\label{sec:rq_discussion}

\noindent\textbf{RQ1. How does joint impairment change gait?}
Hip studies identify changes in hip excursion and compensatory movement of the pelvis and trunk \cite{eitzen2012sagittal,reininga2012subjects}. Knee studies link changes in flexion, step regularity and loading to factors that include pain, speed and body mass \cite{kaufman2001gait,harding2012body,bacon2022relation}. Ankle studies draw attention to dorsiflexion, plantar flexion and propulsion, including residual differences after joint replacement \cite{valderrabano2007gait,schmitt2015effect}. These changes can extend to adjacent joints. The greater coverage of knee motion than ankle motion in Table~\ref{tableCV} also shows why the strength of the evidence varies by joint and plane.

\noindent\textbf{RQ2. What can vision methods recover from walking video?}
Pose estimation provides landmarks that can be tracked over a stride. Gait analysis uses them to estimate variables such as step timing and joint angles, and action recognition uses temporal patterns to assign movement labels \cite{stenum2021two,lv2022analysis,rohan2020human}. The intended output determines the choice of method. Detecting knee hyperextension as a pattern and measuring its peak angle call for different evaluations. Estimates of loading require a biomechanical or learned model whose output is checked against a suitable reference \cite{boswell2021neural}.

\Needspace{6\baselineskip}
\noindent\textbf{RQ3. What evidence supports measurement, severity and change?}
Agreement with a synchronised reference helps establish measurement accuracy when the anatomical definition, plane and gait phase match. A severity claim needs a clinically defined target and evaluation in patients who were held out during model development. Tracking recovery also requires repeatability and evidence that the observed change exceeds measurement error and relates to clinical function \cite{stenum2024clinical,schoenwether2025reliability}. Table~\ref{tab:validation_claims} sets out these requirements. The fall-related studies in Table~\ref{tab:fall_evidence} support recognition of scripted patterns, while a future-fall claim requires prospective outcomes.

\noindent\textbf{RQ4. When do recording conditions and the study population limit use?}
Camera placement and limb visibility affect the accuracy of joint angles. Occluding the far leg, for instance, can widen the limits of agreement during walking \cite{wade2023examination}. Clothing and lighting also affect landmark tracking, and validation across everyday conditions remains uneven (Section~\ref{sec:recording_robustness}). Results from healthy volunteers establish performance for the participants and conditions studied. Evaluating new patients, recording setups and sites helps determine where a method can be used and how often it produces a reportable measurement.
